\chapter{Bewertung und Ergebnis}
\label{cha:Bewertung und Ergebnis}

\section{Bewertung der Testergebnisse}
Die durchgeführten Funktionstests zeigen, dass die wesentlichen Anforderungen an die entwickelte Applikation erfüllt wurden. Die nicht vollständig erfolgreiche Validierung der vollständigen Ablauferkennung ist auf Einschränkung der Testumgebung und auf das Verbinden während des Betriebs zurückzuführen. Durch eine fehlende Unterstützung aller Automatisierungsabläufe der Testlinie, konnten nicht alle in der Applikation implementierte Abläufe getestet werden. Die fehlerhafte Verarbeitung der Prozessabläufe, die als erstes erkannt werden, kann darauf zurückgeführt werden, dass die Applikation eine Grundlage an Informationen für korrekte Erkennung von Prozessen benötigt. Diese Informationen werden jedoch erst im Verlauf der aktiven Verbindung über MQTT geschickt, wodurch Abläufe dann falsch erkannt werden. Dieses Fehlverhalten der Anwendung könnte durch Hinzufügen von Bedingungen behoben werden, wodurch Abläufe erst erkannt werden sollen, wenn genügend Informationen bereitgestellt sind. Da sich aber die sofortige Aufzeichnung von Automatisierungsabläufen als wichtiger eingestuft wird, als jeden einzelnen Prozess mit der richtigen Information abzuschließen, wird diese Einschränkung als tolerierbar bewertet.

Die definierten Performance-Ziele konnten nicht vollständig erreicht werden. Insbesondere die CPU-Auslastung überschritt die festgelegten Grenzwerte während mehrerer Testläufe. Ursache hierfür war sowohl die bewusst hohe und kontinuierliche Nachrichtenlast während der Messungen als auch die rechenintensiven Einstellungen mit vier Nachrichtenfenstern auf einer Seite sowie das Laden von 60 Nachrichten pro Nachrichtenfenster. Durch das verringern von den zu ladenden Nachrichten und gleichzeitig dargestellten Nachrichtenfenstern kann die CPU-Auslastung drastisch reduziert werden. Die Anwendung blieb trotz dieser rechenintensiven Einstellungen jederzeit stabil und konnte alle Nachrichten ohne Datenverluste verarbeiten. 

Neben der überschrittenen CPU-Auslastung ist auch der Arbeitsspeicherverbrauch zwei- bis dreimal höher als der geforderte Grenzwert. Auch bei mehreren Tests pendelt sich der Verbrauch wie bei Abbildung \ref{fig:Messwerte} ungefähr bei 1200 Megabyte ein. Diese Phänomen ist dem sogenannten \textit{Garbagecollector} des \textit{.Net}-Frameworks geschuldet. Dieser ist für das löschen von nicht mehr referenzierten Objekten und toten Instanzen von Referenztypen zuständig, die nicht mehr gebraucht werden, aber Arbeitsspeicher verbrauchen \autocite[22]{garbagecollector}. Da der Garbage Collector aber erst von gewissen Triggern ausgelöst wird wie beispielsweise Arbeitsspeichermangel, arbeitet dieser nicht ständig \autocite[372--375]{garbagecollector}. Das ständige löschen von Objekten würde außerdem auch zu Performance-Einbußen führen \autocite[372]{garbagecollector}. Für ein tieferes Verständnis für den Garbage Collector kann \autocite{garbagecollector} herangezogen werden.

Der Arbeitsspeicherverbrauch der Applikation wird also automatisch durch das \textit{.Net}-Framework organisiert \autocite[135]{garbagecollector}. Aus diesem Grund kann die Performance trotz fehlgeschlagener Tests als ausreichend bewertet werden.
\section{Umsetzung der Anforderungen}
Die Anforderung zur Anzeige, Speicherung und Filterung von MQTT-Nachrichten wurde vollständig umgesetzt. Die Applikation ermöglicht sowohl die Filterung nach Topics als auch die Volltextsuche innerhalb der gespeicherten Nachrichten. Darüber hinaus können drei verschiedene Anzeigemöglichkeiten der Nachrichten ausgewählt werden, wodurch die Anwendung personalisierbar wird.

Die automatische Erkennung von Strukturfehlern sowie verschiedener Prozessfehler wurde erfolgreich implementiert und konnte durch die Funktionstests validiert werden. Neben der reinen Anzeige von Fehlern können diese nach Ablauf- beziehungsweise Strukturfehlern gefiltert werden. Außerdem wurde das Paging-Konzept auch für die Fehler umgesetzt, wodurch diese ressourcenarm angezeigt werden können.

Die Teilnehmerüberwachung erkennt automatisch neue Teilnehmer und visualisiert deren Zustände sowie deren Verbindungsstatus, wodurch die Anforderungen dafür erfüllt wurden.

Bei der Ablauferkennung bildeten sich bei der Durchführung der Funktionstests geringfügige Fehler heraus. Trotz dieser Einschränkungen bietet die implementierte Ablauferkennung einen deutlichen Mehrwert für die Analyse von Produktionsabläufen. Prozessereignisse können automatisch erkannt und fehlerhafte Abläufe durch entsprechende Fehlermeldungen hervorgehoben werden. Dadurch wird die Fehlersuche gegenüber einer rein manuellen Analyse erheblich vereinfacht.

Durch die Nutzung asynchroner Verarbeitung sowie eines Paging-Konzeptes konnte eine stabile Verarbeitung großer Nachrichtenmengen erreicht werden. Die definierten Zielwerte für CPU- und Arbeitsspeicherverbrauch wurden jedoch nicht in allen Testfällen eingehalten. 

Durch die Verwendung der MVVM-Architektur sowie der Trennung von Services, Repositories und Benutzeroberfläche wurde eine Grundlage für zukünftige Erweiterungen geschaffen. Für den Fall eines zukünftigen Umstiegs auf andere Betriebssysteme ist die Applikation nicht sofort übertragbar und einsetzbar. Dafür müsste die Anwendung auf ein anderes Framework übertragen werden, was jedoch durch die MVVM-Architektur stark vereinfacht wird.

Die Anforderung an konfigurierbare Prüfmechanismen erfüllt der MQTT Analyzer nicht. Dafür ist der Prozessablauf einer jeden Maschine zu modular aufgebaut, wodurch die Umsetzung dieser Anforderung für den Rahmen der vorliegenden Arbeit sprengen würde.

Insgesamt zeigt die Bewertung, dass die entwickelte Anwendung für den Einsatz im Produktionsumfeld geeignet ist. Durch die Kombination aus Nachrichtenanalyse, automatischer Fehlererkennung, Teilnehmerüberwachung und Ablauferkennung bietet sie einen deutlichen Mehrwert gegenüber den bisher eingesetzten MQTT-Werkzeugen. Der Anwender wird dadurch bei der Inbetriebnahme, Analyse und Fehlerdiagnose von Verpackungslinien wirksam unterstützt.

\section{Vergleich mit bestehenden Werkzeugen}
Die in Kapitel \ref{cha:Ausgangssituation} beschriebenen Werkzeuge bieten nur eingeschränkte Möglichkeiten zur Analyse von MQTT-Nachrichten. Die meisten Werkzeuge sind auf die reine Anzeige von Nachrichten beschränkt und bieten keine erweiterten Funktionen wie Filterung, Volltextsuche oder automatische Fehlererkennung. Der MQTT Analyzer hingegen bietet eine umfassende Lösung für die Analyse von MQTT-Nachrichten und unterstützt den Anwender bei der Inbetriebnahme, Analyse und Fehlerdiagnose von Verpackungslinien. Für einen genaueren Vergleich zeig Tabelle \ref{tab:Werkzeugvergleich} die wesentlichen Unterschiede und Vorteile der einzelnen Applikationen. 

\begin{table}[htbp] 
    \centering \caption{Nutzwertanalyse der untersuchten MQTT-Werkzeuge} 
    \label{tab:Werkzeugvergleich} 
    \begin{tabular}{lccccc} \hline \textbf{Kriterium} & \textbf{MQTT.fx} & \textbf{Explorer} & \textbf{Multimeter} & \textbf{Analyzer} \\ 
        \hline Nachrichtenanzeige (5\%) & 3 & 4 & 4 & 5 \\ 
        Topic-Filterung (15\%) & 3 & 4 & 3 & 5 \\ 
        Volltextsuche (15\%) & 1 & 5 & 3 & 5 \\ 
        Persistente Speicherung (10\%) & 1 & 1 & 2 & 5 \\ 
        Analyse historischer \-Daten (10\%) & 3 & 2 & 1 & 5 \\ 
        Teilnehmerüberwachung (10\%) & 1 & 1 & 1 & 5 \\ 
        Automatische Fehler\-erkennung (10\%) & 1 & 1 & 1 & 5 \\ 
        Automatisierungsablauf\-analyse (10\%) & 1 & 1 & 1 & 3 \\ 
        Nachrichtensimulation (5\%) & 3 & 1 & 2 & 3 \\ 
        Erweiterbarkeit (10\%) & 1 & 1 & 3 & 4 \\ 
        \hline Gesamtpunktzahl & 1.7 & 2.3 & 2.1 & 4.6 \\ 
        \hline 
    \end{tabular}
\end{table}

Die Tabelle zeigt, dass der MQTT Analyzer in allen relevanten Bereichen überlegen ist und somit eine wertvolle Erweiterung für die zur Verfügung stehenden Analyse- und Diagnoseanwendungen ist. Vor allem durch die Fehler- und Ablauferkennung sowie Filterung der Nachrichten bietet die selbst entwickelte Anwendung eine Möglichkeit, Fehler schneller zu erkennen und Analysezeiten zu reduzieren. 

Die Zeitersparnis wirkt sich unmittelbar auf die Wirtschaftlichkeit aus. Kürzere Analyse- und Diagnosezeiten führen zu einem geringeren Personalaufwand während der Inbetriebnahme und Fehlerbehebungen von Verpackungslinien. Gleichzeitig können Kommunikationsprobleme schneller identifiziert und behoben werden, wodurch Stillstände von Maschinen reduziert werden können.